\section{Обсуждение: условия выбора и границы результата}
\paragraph{Что следует учитывать при выборе.}
Проверенные эффекты поддерживают выбор по сочетанию трёх характеристик:
механизма формирования парной разметки, реализованной параметризации
полезности и эксплуатационной квоты. Сравнительная деградация
gain-методов против Weak-error повторяется на втором датасете;
внутри gain-семейства найдены конкретные преимущества и взаимодействия.
Корректировка известного отбора даёт отдельный измеримый прирост,
поэтому raw-победитель и победитель среди корректно обученных
кандидатов отвечают на разные вопросы.

Практически это означает: сохранять простой обучаемый baseline,
проверять попарные взаимодействия с механизмом отбора, сравнивать
IPW и raw на одинаковых извлечениях и оценивать окончательный
выбор на отдельном dev при нужной квоте. Полученные условия
предпочтительности являются локальными ориентирами для такой процедуры,
а не правилом всегда выбирать Weak-error, Hurdle или Ternary.

\paragraph{Как интерпретировать неподтверждённые сравнения.}
Полная проверка сохраняет важные границы: в семье B natural-различия
и gain-only взаимодействия не прошли поправку; в C не установлены
прямое преимущество accuracy-настройки над MSE или надёжная смена
знака преимущества между natural и harm8. В семье M найден один
контраст методов по accuracy и macro-F1; случаев различимой macro-P/F1
при необнаруженной accuracy-разнице в той же семье нет.
Эти результаты ограничивают область утверждений, но не отменяют
конкретных эффектов таблицы~\ref{tab:main_effects}.
Незначимость не доказывает эквивалентности; разные статусы значимости
не заменяют прямой тест различия эффектов.

\paragraph{Функциональный класс и критерий настройки.}
Direct --- линейный ridge-прогноз; Hurdle --- произведение двух
logistic-голов. Поэтому найденные различия относятся к полным
реализованным способам обучения и могут включать вклад ёмкости,
потерь и регуляризации. Accuracy-настройка меняет выбранную
регуляризацию примерно в 38--48\% случаев в рассмотренных основных
срезах, но сама смена tuning-критерия не получила подтверждённого
преимущества после общей поправки. Восемь gain-only взаимодействий
относятся именно к заданному accuracy-tuned протоколу;
их наличие не является доказательством превосходства этого критерия.

\paragraph{Обобщение и статистический статус.}
BANKING77 повторно использовался после просмотра исходных результатов;
CLINC150 добавлен по отдельному плану, но его multiclass-разбор
выполнен уже после основных экспериментов. Выбранные иллюстрации
не образуют новую подтверждающую выборку. Панели внутри датасетов
зависимы, bootstrap условен на пулах, экспертах и обученных scores,
а два intent-датасета не представляют всё разнообразие задач routing.
Квота измеряет число вторых вызовов, не фактические денежные затраты
или end-to-end latency. Отбор по парным исходам предполагает,
что разметка полного пула уже приобретена.

\paragraph{Следующий содержательный контроль.}
Для отделения обучающей цели от функционального класса нужен
одинаковый Hurdle-подобный нелинейный score, обученный либо прямой MSE
на приросте, либо факторизованными потерями, с одинаковыми признаками,
регуляризацией и tuning-бюджетом. Затем критерий выбора по условной
структуре ошибок и устойчивости границы квоты следует зафиксировать
на development-данных и проверить на новой задаче. Этот контроль
ещё не выполнялся; настоящие результаты задают его мотивацию
и измеримые эффекты, которые требуется объяснить.
